\scanpage{35}
\mspage{17}
inverse de $x$ dans le ``rév. universel abélien'' $\widetilde G$ de $G$, construit comme $\varprojlim$ des revêtements $G(n)\simeq G$ de $G$, donnés par $n\operatorname{id}$. L'énoncé dit que si ce torseur est trivial -- i.e. si $\forall n$, $x$ se remonte -- alors $x=0$. C'est vrai en car. $p>0$ également -- mais dans ce cas les ${}_nG$ ne sont pas étales.

\margin{\deleted{\unsure{Le théorème} \illegible{2} \unsure{Attention} \illegible{2} \unsure{de cette question} l'\unsure{affirmation} \illegible{3} \unsure{fausse}...}}

Si cependant on remplace les ${}_{p^\nu}G$ par leurs quotients par leurs parties radicielles -- en travaillant donc dans la cat. des ``quasi-variétés'' où on \deleted{\unsure{identifie}} \inserted{inverse} les \deleted{\unsure{immersions|isomorphismes}} \inserted{morphismes} finis radiciels surjectif\unsure{s}, \deleted{\unsure{je pense qu'il}} \inserted{pour travailler} avec la partie proétale \unsure{du|de} $H_{\unsure{1|n}}(\overline G)$, \deleted{\unsure{qu'il en|quitte}} l'existence \unsure{des|de} \unsure{relèvts|relèvements} radiciels $/K$ implique déjà $x=0$. \inserted{Bien sur $x$ est de $p$-torsion,} On est ramené au cas d'une VA, avec $px=0$.... Mais \deleted{la} validité sous cette forme du th. 3 ou \inserted{de son cor. en} car. $p>0$ nous fait une belle jambe -- car pour être à \unsure{l'aise|l'air}, il faudrait pouvoir les prouver \deleted{pour} \inserted{dans} les contextes des groupes fondamentaux ``premiers à $p$''. Or dans ce con\inserted{texte''}, déjà dans le cas d'une courbe elliptique $Y$ (\deleted{\unsure{courbe}} \inserted{\unsure{rendue}} anabélienne au besoin en enlevant des pts) l'injectivité de (**) est \uline{fausse}. Il

\margin{O.K. en se remplaçant d'abord $K$ par \unsure{un|une} \unsure{alg|alg.} de type fini sur $\mathbb F_p$, puis en se ramenant \unsure{au cas} d'un corps fini \uline{parfait}...}
\margin{34}
